\documentclass[english, parskip=full]{scrartcl}
\usepackage[utf8]{inputenc}  % Kodierung der Datei
\usepackage[ngerman]{babel}  % Sprache auf Deutsch 
\usepackage{color}
\usepackage{graphicx, gensymb}
\usepackage{amsfonts, cancel, amssymb}
\usepackage{amsmath, amsthm, array, esint}
\usepackage{booktabs, multirow}  % schönere Tabellen
\usepackage{caption}  % Anpassung der Captions
\usepackage{scrlayer-scrpage}    % Manipulation des Seitenstils
\pagestyle{scrheadings}  % Stil für die Seite setzen
\automark{section}  % Abschnittsnamen als Seitenbeschriftung 
\setheadsepline{.5pt}    % Kopzeile durch Linie abtrennen
\usepackage[hidelinks]{hyperref}  % Links und weitere PDF-Features
\usepackage{typearea, tikz, pgffor, verbatim}
\usetikzlibrary{calc, arrows.meta,arrows, graphs, quotes, positioning}
\usepackage[cache=false, outputdir=build]{minted}
\usemintedstyle{vs}
\usepackage{placeins} % provides \FloatBarrier

\definecolor{bg}{rgb}{0.9,0.9,0.9}
\newcommand\numberthis{\addtocounter{equation}{1}\tag{\theequation}}
\usepackage{svg}
\usepackage{siunitx}
\usepackage{physics}
\usepackage{tensor}
\renewcommand{\bra}{}
\newcommand{\kom}{}
\newcommand{\brak}{}
\renewcommand{\braket}{}
\newcommand{\intx}{}
\newcommand{\intr}{}
\newcommand{\kla}{}
\newcommand{\klam}{}
\newcommand{\klammer}{}
\newcommand{\oper}{}
\newcommand{\tecxt}{}
\newcommand{\Bra}{}
\newcommand{\Ket}{}
\renewcommand{\i}{\text{i}}
\newcommand{\e}{\text{e}}
\newcommand*{\Fourier}{\textsc{Fourier}\ }
\newcommand*{\F}{\mathcal{F}}
\newcommand*{\sinc}{\text{sinc\,}}
\newcommand{\conv}{\mathop{\scalebox{3.5}{\raisebox{-0.38ex}{\makebox[0.5\width][c]{$\ast$}}}}\limits}

\newcommand*{\titel}{Phononen in zweiter Quantisierung}
\newcommand*{\autor}{Joachim Schwardt}

\hypersetup{pdfauthor={\autor}, pdftitle={\titel}}

%\titlehead{titlehead}
\subject{Theoretische Festkörperphysik}
\title{\titel}
\author{\autor}
\date{\today}

\begin{document}

%\begin{titlepage}
\maketitle  % Titel setzen
%\tableofcontents  % Inhaltsverzeichnis setzen
%\end{titlepage}

\section{Isotropie hexagonaler Kristalle}
Für einen hexagonalen Kristall hat der elastische Tensor in Voigt'scher Notation die Form
\begin{align}
C &= \begin{pmatrix}
C_{11} & C_{12} & C_{13} & 0 & 0 & 0 \\
C_{12} & C_{11} & C_{13} & 0 & 0 & 0 \\
C_{13} & C_{13} & C_{33} & 0 & 0 & 0 \\
0 & 0 & 0 & C_{44} & 0 & 0 \\
0 & 0 & 0 & 0 & C_{44} & 0 \\
0 & 0 & 0 & 0 & 0 & C_{66} \\
\end{pmatrix},
\end{align}
wobei $C_{66} = \frac{C_{11} - C_{12}}{2}$.
Die Indizes stehen für $1\equiv xx, 2\equiv yy, 3\equiv zz, 4\equiv yz, 5\equiv xz, 6\equiv xy$.
Die Christoffel-Gleichungen lauten
\begin{align}
0 &= \sum_{j'} \klam\left[ \tilde{D}_{jj'} - \rho c^2(\hat{q})\delta_{jj'} \right]e_{j'}(\vec{q}) && \tecxt\text{mit} & \tilde{D}_{jj'} &= \sum_{kk'} C_{jk,j'k'} \hat{q}_k \hat{q}_{k'}. 
\end{align}
Wir finden die sechs Komponenten
\begin{align}
\tilde{D}_{xx} &= C_{xx,xx}\hat{q}_x\hat{q}_x + C_{xy,xy}\hat{q}_y\hat{q}_y + C_{xz,xz}\hat{q}_z\hat{q}_z,\\
\tilde{D}_{xy} &= C_{xx,yy}\hat{q}_x\hat{q}_y + C_{xy,yx}\hat{q}_y\hat{q}_x,\\
\tilde{D}_{xz} &= C_{xx,zz}\hat{q}_x\hat{q}_z + C_{xz,zx}\hat{q}_z\hat{q}_x,\\
\tilde{D}_{yy} &= C_{yx,yx}\hat{q}_x\hat{q}_x + C_{yy,yy}\hat{q}_y\hat{q}_y + C_{yz,yz}\hat{q}_z\hat{q}_z,\\
\tilde{D}_{yz} &= C_{yy,zz}\hat{q}_y\hat{q}_z + C_{yz,zy}\hat{q}_z\hat{q}_y,\\
\tilde{D}_{zz} &= C_{zx,zx}\hat{q}_x\hat{q}_x + C_{zy,zy}\hat{q}_y\hat{q}_y + C_{zz,zz}\hat{q}_z\hat{q}_z.
\end{align}
In Matrixschreibweise (symmetrische Einträge ausgelassen)
\begin{align}
\tilde{D} &= \begin{pmatrix}
C_{11}\hat{q}_x^2 + C_{66}\hat{q}_y^2 + C_{44}\hat{q}_z^2 & (C_{12} + C_{66})\hat{q}_x\hat{q}_y & (C_{13} + C_{44}) \hat{q}_x\hat{q}_z \\
& C_{66}\hat{q}_x^2 + C_{11} \hat{q}_y^2 + C_{44}\hat{q}_z^2 & (C_{13} + C_{44}) \hat{q}_y\hat{q}_z \\
& & C_{44}\hat{q}_x^2 + C_{44}\hat{q}_y^2 + C_{33} \hat{q}_z^2
\end{pmatrix}.
\end{align}
Mit den Komponenten $e_{x,y,z}$ der Polarisationsvektoren $\vec{e}(\vec{q})$ lauten die Christoffel-Gleichungen
\begin{align}
\rho c^2(\hat{q}) e_x &= \kla\left( C_{11}\hat{q}_x^2 + C_{66}\hat{q}_y^2 + C_{44}\hat{q}_z^2 \right)e_x + (C_{12} + C_{66})\hat{q}_x\hat{q}_y e_y + (C_{13} + C_{44}) \hat{q}_x\hat{q}_z e_z,\\
\rho c^2(\hat{q}) e_y &= (C_{12} + C_{66})\hat{q}_x\hat{q}_y e_x + \kla\left( C_{66}\hat{q}_x^2 + C_{11} \hat{q}_y^2 + C_{44}\hat{q}_z^2 \right)e_y + (C_{13} + C_{44}) \hat{q}_y\hat{q}_z e_z,\\
\rho c^2(\hat{q}) e_z &= (C_{13} + C_{44}) \hat{q}_x\hat{q}_z e_x + (C_{13} + C_{44}) \hat{q}_y\hat{q}_z e_y + \kla\left( C_{44}\hat{q}_x^2 + C_{44}\hat{q}_y^2 + C_{33} \hat{q}_z^2 \right) e_z.
\end{align}
Für $\hat{q}_z=0$ vereinfachen sich die Gleichungen zu
\begin{align}
\rho c^2(\hat{q}) e_x &= \kla\left( C_{11}\hat{q}_x^2 + C_{66}\hat{q}_y^2 \right)e_x + (C_{12} + C_{66})\hat{q}_x\hat{q}_y e_y,\\
\rho c^2(\hat{q}) e_y &= (C_{12} + C_{66})\hat{q}_x\hat{q}_y e_x + \kla\left( C_{66}\hat{q}_x^2 + C_{11} \hat{q}_y^2 \right)e_y,\\
\rho c^2(\hat{q}) e_z &= C_{44} \kla\left( \hat{q}_x^2 + \hat{q}_y^2 \right) e_z = C_{44} e_z.
\end{align}
Die letzte Gleichung liefert $e_z=0$ oder $c^2(\hat{q}) = C_{44}$.
Da wir aber an der Phasengeschwindigkeit entlang der Basisfläche interessiert sind, wählen wir hier $e_z=0$.
Umstellen der zweiten Gleichung ergibt
\begin{align}
e_y &= \frac{C_{12} + C_{66}}{\rho c^2(\hat{q}) - C_{66}\hat{q}_x^2 - C_{11} \hat{q}_y^2}\hat{q}_x\hat{q}_y e_x.
\end{align}
Einsetzen in die erste Gleichung ergibt
\begin{align*}
\rho c^2(\hat{q}) - C_{11}\hat{q}_x^2 - C_{66} \hat{q}_y^2 &= (C_{12} + C_{66}) \hat{q}_x \hat{q}_y \frac{C_{12} + C_{66}}{\rho c^2(\hat{q}) - C_{66}\hat{q}_x^2 - C_{11} \hat{q}_y^2}\hat{q}_x\hat{q}_y, \numberthis \\
\Rightarrow\ \ \ (C_{12} + C_{66})^2\hat{q}_x^2 \hat{q}_y^2 &= \rho^2 c^4(\hat{q}) + C_{11}C_{66} (\hat{q}_x^4 + \hat{q}_y^4) - \\
&\hspace{0.5cm} - \rho c^2(\hat{q}) \kla\left( C_{11} + C_{66} \right) (\hat{q}_x^2 + \hat{q}_y^2) + (C_{11}^2 + C_{66}^2) \hat{q}_x^2\hat{q}_y^2 . \numberthis
\end{align*}
Das ist eine quadratische Gleichung in $\rho c^2(\hat{q})$, wobei $\hat{q}_x^2 + \hat{q}_y^2 = 1$ und
\begin{align}
(C_{12} + C_{66})^2 - C_{11}^2 - C_{66}^2 &= \frac{(C_{11} + C_{12})^2}{4} - C_{11}^2 - \frac{(C_{11} - C_{12})^2}{4} = \\
&= C_{11}C_{12} - C_{11}^2 = C_{11}(C_{12} - C_{11}) = -2C_{11}C_{66}
\end{align}
gelten.
Wir lösen also die Gleichung
\begin{align}
0 &= \rho^2 c^4(\hat{q}) - \rho c^2(\hat{q}) (C_{11} + C_{66}) + C_{11}C_{66} (\hat{q}_x^4 + \hat{q}_y^4) + 2C_{11}C_{66} \hat{q}_x^2 \hat{q}_y^2,\\
\Rightarrow\ \ \ \rho c^2(\hat{q}) &= \frac{C_{11} + C_{66}}{2} \pm \sqrt{\frac{(C_{11} + C_{66})^2}{4} - C_{11}C_{66} \kla\left(  \hat{q}_x^2 + \hat{q}_y^2 \right)^2} = \\
&= \frac{C_{11} + C_{66}}{2} \pm \sqrt{\frac{(C_{11} - C_{66})^2}{4}} = \begin{cases} C_{11} \\ C_{66} \end{cases}.
\end{align}
Im letzten Schritt haben wir $C_{11} \ge C_{66}$ angenommen, andernfalls ist die Fallunterscheidung gerade andersherum.
Unabhängig davon ist die Phasengeschwindigkeit aber isotrop entlang der Basisfläche, weil wir keine Abhängigkeit von $\hat{q}$ haben.
\\
Jetzt können wir die Polairsationsvektoren einfach über den Zusammenhang von $e_x$ und $e_y$ bestimmen
\begin{align}
\vec{e}(\vec{q}) &= \frac{1}{\mathcal{N}}\begin{pmatrix}
\rho c^2 - C_{66} \hat{q}_x^2 - C_{11} \hat{q}_y^2 \\ 
(C_{12} + C_{66} ) \hat{q}_x \hat{q}_y \\ 
0
\end{pmatrix} = \frac{1}{\mathcal{N}}\begin{pmatrix}
\rho c^2 - C_{66} \hat{q}_x^2 - C_{11} \hat{q}_y^2 \\ 
(C_{11} - C_{66} ) \hat{q}_x \hat{q}_y \\ 
0
\end{pmatrix}.
\end{align}
mit einer zunächst nicht näher bestimmten Normierung.
Für $\rho c^2 = C_{11}$ ergibt sich
\begin{align}
\vec{e}_1(\vec{q}) &= \frac{(C_{11} - C_{66})}{\mathcal{N}} \begin{pmatrix}
\hat{q}_x^2 \\ \hat{q}_x\hat{q}_y \\ 0
\end{pmatrix} \equiv \vec{q},
\end{align}
und analog für $\rho c^2 = C_{66}$ 
\begin{align}
\vec{e}_2(\vec{q}) &= \frac{(C_{11} - C_{66})}{\mathcal{N}} \begin{pmatrix}
-\hat{q}_y^2 \\ \hat{q}_x\hat{q}_y \\ 0
\end{pmatrix} \equiv \begin{pmatrix}
-\hat{q}_y \\ \hat{q}_x \\ 0
\end{pmatrix}.
\end{align}
%Die Normierungskonstante ist
%\begin{align}
%\mathcal{N}^2 &= c^4 + C_{66}^2 \hat{q}_x^4 + C_{11}^2 \hat{q}_y^4 - 2C_{11}C_{66} \hat{q}_x^2 \hat{q}_y^2 - 2c^2 (C_{11}\hat{q}_x^2 + C_{66}\hat{q}_y^2) + (C_{11} - C_{66})^2 \hat{q}_x^2\hat{q}_y^2 = \\
%&= c^4
%\end{align}

\section{Harmonische Näherung in zweiter Quantisierung}
Gegeben sei ein $d$-dimensionaler kristall mit $N$ Gitterplätzen und einem Atom der Masse $M$ pro Elementarzelle.
In harmonischer Näherung gilt
\begin{align}
\mathcal{H}_\tecxt\text{h} &= \sum_{\textbf{q} \lambda} \hbar\omega_\lambda (\textbf{q}) \kla\left( b^\dagger_{\textbf{q} \lambda} b_{\textbf{q} \lambda} + \frac{1}{2} \right).
\end{align}
Die Operatoren $\textbf{u}_n$ für die Auslenkung der Gitteratome lassen sich durch Phononen-Erzeugungs- und Vernichtungsoperatoren $b_{\textbf{q} \lambda}^\dagger$ und $b_{\textbf{q} \lambda}$ ausdrücken
\begin{align}
\textbf{u}_n &= \frac{1}{\sqrt{NM}} \sum_{\textbf{q} \in \tecxt\text{1.BZ}, \lambda} \sqrt{\frac{\hbar}{2\omega_\lambda(\textbf{q})}}  \e^{\i \textbf{q}\cdot \textbf{R}_n^0} \textbf{e}_{\textbf{q} \lambda} \kla\left( b_{\textbf{q} \lambda} + b^\dagger_{-\textbf{q} \lambda} \right).
\end{align}
\subsection*{(a) Mittlere thermische Auslenkung}
Die Zustandssumme für diesen Hamiltonian ist
\begin{align}
Z &= \tecxt\text{tr\,} \e^{-\beta \mathcal{H}_\tecxt\text{h}} = \sum_{\{n_{\textbf{q} \lambda}\}} \langle \{ n_{\textbf{q} \lambda}\} | \prod_{\textbf{q} \lambda} \e^{-\beta\hbar\omega_\lambda (\textbf{q}) \kla\left( b^\dagger_{\textbf{q} \lambda} b_{\textbf{q} \lambda} + \frac{1}{2} \right)} | \{n_{\textbf{q} \lambda}\} \rangle = \\
&= \prod_{\textbf{q} \lambda} \sum_{n_{\textbf{q} \lambda} = 0}^\infty  \e^{-\beta\hbar\omega_\lambda (\textbf{q}) \kla\left( n_{\textbf{q} \lambda} + \frac{1}{2} \right)} = \\
&= \prod_{\textbf{q} \lambda} \frac{\e^{-\beta\hbar\omega_\lambda (\textbf{q}) / 2}}{1 - \e^{-\beta\hbar\omega_\lambda (\textbf{q})}} = \prod_{\textbf{q} \lambda} \frac{1}{2} \csch\kla\left( \frac{\beta\hbar\omega_\lambda (\textbf{q})}{2} \right).
\end{align}
Erwartungswerte können wir dann über $\bra\langle A \rangle = \tecxt\text{tr\,}A\hat{\rho}$ bestimmen.
Hier benötigen wir
\begin{align}
\bra\langle b_{\textbf{q} \lambda} \rangle &= \sum_{\{n_{\textbf{q} \lambda}\}} \langle \{ n_{\textbf{q} \lambda}\} | b_{\textbf{q} \lambda} \hat{\rho} | \{n_{\textbf{q} \lambda}\} \rangle = \\
&= \sum_{\{n_{\textbf{q} \lambda}\}} \langle \dots, n_{\textbf{q} \lambda} + 1, \dots | \sqrt{n_{\textbf{q} \lambda} + 1} \hat{\rho} | \{n_{\textbf{q} \lambda}\} \rangle = 0.
\end{align}
Daraus folgt für die mittlere thermische Auslekung dann
\begin{align}
\bra\langle \textbf{u}_n \rangle &= \dots \bra\langle b_{\textbf{q} \lambda} + b^\dagger_{-\textbf{q} \lambda} \rangle = 0.
\end{align}
\subsection*{(b) Mittlere Abweichung der Auslenkung}
%Wir berechnen zunächst das Quadrat der Auslenkung, wobei wir $\textbf{e}_{-\textbf{q}\lambda} = \textbf{e}_{\textbf{q} \lambda}^\ast \equiv \textbf{e}_{\textbf{q} \lambda}$ und $\textbf{e}_{\textbf{q} \lambda} \cdot \textbf{e}_{\textbf{q} \mu} = \delta_{\lambda\mu}$ ausnutzen.
%Die erste Relation gilt, weil wir eine einatomige Basis haben, weshalb die Polaritsationsvektoren reell gewählt werden können.
Um die Erwartungswerte der Auslenkungsquadrate auszurechnen, nutzen wir $\textbf{e}_{-\textbf{q}\lambda} = \textbf{e}_{\textbf{q} \lambda}^\ast \equiv \textbf{e}_{\textbf{q} \lambda}$ und $\textbf{e}_{\textbf{q} \lambda} \cdot \textbf{e}_{\textbf{q} \mu} = \delta_{\lambda\mu}$.
Die erste Relation gilt, weil wir eine einatomige Basis haben, weshalb die Polaritsationsvektoren reell gewählt werden können.
%Insbesondere folgt dann mit $\textbf{u}_n^\dagger = \textbf{u}_n$, dass das Vorzeichen im Exponenten des Terms $\e^{\i \textbf{q}\cdot \textbf{R}_n^0}$ irrelevant ist.
Damit lauten die Auslenkungsquadrate dann zunächst
\begin{align}
\textbf{u}_n^2 &= \frac{\hbar}{2NM} \sum_{\textbf{q}\textbf{q}' \lambda \lambda'} \frac{\textbf{e}_{\textbf{q} \lambda} \cdot \textbf{e}_{\textbf{q}' \lambda'} \e^{\i (\textbf{q} + \textbf{q}') \cdot \textbf{R}_n^0} }{\sqrt{\omega_\lambda(\textbf{q}) \omega_{\lambda'} (\textbf{q}')}}  \kla\left( b_{\textbf{q} \lambda} b_{\textbf{q}' \lambda'} + b_{\textbf{q} \lambda} b^\dagger_{-\textbf{q}' \lambda'} + b^\dagger_{-\textbf{q} \lambda} b_{\textbf{q}' \lambda'} + b^\dagger_{-\textbf{q} \lambda} b^\dagger_{-\textbf{q}' \lambda'} \right).
\end{align}
Die Erwartungswerte der Erzeuge und Vernichter vereinfachen sich zu
\begin{align}
\bra\langle b_{\textbf{q}\lambda}b_{\textbf{q}'\lambda'} &+ b_{\textbf{q}\lambda} b_{-\textbf{q}'\lambda'}^\dagger + b_{-\textbf{q}\lambda}^\dagger b_{\textbf{q}'\lambda'} + b_{-\textbf{q}\lambda}^\dagger b_{-\textbf{q}'\lambda'}^\dagger \rangle = \\
&= \sum_{\{n_{\textbf{k}\mu} \}} \langle \{n_{\textbf{k}\mu} \} | \left(  b_{\textbf{q}\lambda} b_{-\textbf{q}'\lambda'}^\dagger + b_{-\textbf{q}\lambda}^\dagger b_{\textbf{q}'\lambda'} \right) \hat{\rho} | \{n_{\textbf{k}\mu}\} \rangle = \\
&= \delta_{\textbf{q},-\textbf{q}'} \delta_{\lambda\lambda'} \sum_{\{n_{\textbf{k}\mu} \}} \langle \{n_{\textbf{k}\mu} \} | \left(  2b_{\textbf{q}\lambda}^\dagger b_{\textbf{q}\lambda} + 1 \right) \hat{\rho} | \{n_{\textbf{k}\mu}\} \rangle
\end{align}
Wir schreiben den Operator $\hat{n}_{\textbf{q}\lambda} + \frac{1}{2}$ im Erwartungswert als Ableitung von 
\begin{align}
\hat{\rho} = \frac{1}{Z} \exp\kla\left( -\sum_{\textbf{q}\lambda} \beta\hbar \omega_\lambda (\textbf{q}) \kla\left( \hat{n}_{\textbf{q}\lambda} + \frac{1}{2} \right)  \right),
\end{align}
womit sich der vorherige Ausdruck dann zu
\begin{align}
\dots &= 2\delta_{\textbf{q},-\textbf{q}'} \delta_{\lambda\lambda'} \frac{1}{Z} \left(-\frac{\partial_{\omega_\lambda (\textbf{q})}}{\hbar\beta} \right) Z = \frac{2}{\hbar\beta} \partial_{\omega_\lambda (\textbf{q})} \sum_{\textbf{k} \mu} \log \sinh\kla\left( \frac{\beta\hbar\omega_\mu (\textbf{k})}{2} \right) = \\
&= \delta_{\textbf{q},-\textbf{q}'} \delta_{\lambda\lambda'} \coth\kla\left( \frac{\beta\hbar\omega_\lambda ( \textbf{q})}{2} \right)
\end{align}
vereinfacht.
%Wir nutzen noch $\frac{1}{N}\sum_{\textbf{k}} \e^{\i (\textbf{q} + \textbf{k}) \cdot \textbf{R}_n^0} = \delta_{\textbf{q}, -\textbf{k}}$.
%Dann gilt
%\begin{align}
%\textbf{u}_n^2 &= \frac{\hbar}{2NM} \sum_{\textbf{q},\textbf{k}, \lambda, \mu} \frac{\textbf{e}_{\textbf{q} \lambda} \cdot \textbf{e}_{\textbf{k} \mu}}{\sqrt{\omega_\lambda(\textbf{q}) \omega_\mu (\textbf{k})}} \e^{\i (\textbf{q} + \textbf{k})\cdot \textbf{R}_n^0}  \kla\left( b_{\textbf{q} \lambda} b_{\textbf{k} \mu} + b_{\textbf{q} \lambda} b^\dagger_{\textbf{k} \mu} + b^\dagger_{\textbf{q} \lambda} b_{\textbf{k} \mu} + b^\dagger_{\textbf{q} \lambda} b^\dagger_{\textbf{k} \mu} \right) = \\
%&= \frac{\hbar}{2NM} \sum_{\textbf{q},\lambda} \frac{1}{\omega_\lambda (\textbf{q})} \kla\left( b_{\textbf{q} \lambda}^2 + b^{\dagger 2}_{\textbf{q} \lambda} + b_{\textbf{q} \lambda} b^\dagger_{\textbf{q} \lambda} + b^\dagger_{\textbf{q} \lambda} b_{\textbf{q} \lambda} \right) = \\
%&= \frac{\hbar}{2NM} \sum_{\textbf{q} \lambda} \frac{1}{\omega_\lambda (\textbf{q}) } \kla\left( b_{\textbf{q} \lambda}^2 + b^{\dagger 2}_{\textbf{q} \lambda} + 2\hat{n}_{\textbf{q} \lambda} + 1 \right).
%\end{align}
%Die quadratischen Terme liefern im Erwartungswert wieder keinen Beitrag.
%Die anderen beiden Terme lassen sich dann zusammen berechnen, da
%\begin{align}
%\bra\langle \hat{n}_{\textbf{q} \lambda} + \frac{1}{2} \rangle &= \sum_{\{n_{\textbf{k} \mu}\}} \braket\langle \{n_{\textbf{k} \mu}\} | \kla\left( \hat{n}_{\textbf{q} \lambda} + \frac{1}{2} \right) \hat{\rho} | \{n_{\textbf{k} \mu}\} \rangle = \\
%&= \frac{1}{Z} \left(-\frac{\partial_{\omega_\lambda (\textbf{q})}}{\hbar\beta} \right) Z = \frac{1}{\hbar\beta} \partial_{\omega_\lambda (\textbf{q})} \sum_{\textbf{k} \mu} \log \sinh\kla\left( \frac{\beta\hbar\omega_\mu (\textbf{k})}{2} \right) = \\
%&= \frac{1}{2} \coth\kla\left( \frac{\beta\hbar\omega_\lambda ( \textbf{q})}{2} \right).
%\end{align}
%Die mittlere quadratische Abweichung der Auslenkung ergibt sich dann zu
%\begin{align}
%\bra\langle \textbf{u}_n^2 \rangle &= \frac{\hbar}{NM} \sum_{\textbf{q} \lambda} \frac{1}{\omega_\lambda (\textbf{q})} \bra\langle \hat{n}_{\textbf{q} \lambda} + \frac{1}{2} \rangle = \\
%&= \frac{\hbar}{2NM} \sum_{\textbf{q} \lambda} \frac{1}{\omega_\lambda( \textbf{q}) } \coth\kla\left( \frac{\beta\hbar\omega_\lambda (\textbf{q})}{2} \right).
%\end{align}
Die mittlere quadratische Abweichung der Auslenkung ergibt sich dann zu
\begin{align}
\bra\langle \textbf{u}_n^2 \rangle &= \frac{\hbar}{2NM} \sum_{\textbf{q} \lambda} \frac{1}{\omega_\lambda( \textbf{q}) } \coth\kla\left( \frac{\beta\hbar\omega_\lambda (\textbf{q})}{2} \right).
\end{align}
Die Phononen-Zustandsdichte ist definiert als
\begin{align}
n(\omega) &:= \frac{1}{N} \sum_{\textbf{q} \lambda} \delta\kla\left( \omega - \omega_\lambda ( \textbf{q}) \right).
\end{align}
Damit lässt sich die obige Gleichung auch als Integral über $\omega$ ausdrücken, da
\begin{align}
\bra\langle \textbf{u}_n^2 \rangle &= \frac{\hbar}{2NM} \sum_{\textbf{q} \lambda} \intx\int_{ }^{ } \delta \kla\left( \omega - \omega_\lambda (\textbf{q}) \right) \frac{1}{\omega} \coth\kla\left( \frac{\beta\hbar\omega}{2} \right) \, \mathrm{d}\omega = \\
&= \frac{\hbar}{2M} \intx\int_{ }^{ }n(\omega) \frac{1}{\omega} \coth\kla\left( \frac{\beta\hbar\omega}{2} \right) \, \mathrm{d}\omega.
\end{align}
\subsection*{(c) Divergenz für $d<3$}
Wir betrachten den Integranden zunächst für kleine $x := \frac{\beta\hbar\omega}{2} \ll 1$.
Dann gilt $n(x) \sim x^{d-1}$ und wir finden die Entwicklung
\begin{align}
\frac{n(x)}{x} \coth(x) &\sim x^{d-2} \frac{1 + \frac{1}{2}x^2}{x + \frac{1}{6}x^3} \sim x^{d-3} \kla\left( 1 + \frac{x^2}{2} \right) \kla\left( 1 - \frac{x^2}{6} + \dots \right) \sim \\
&\sim x^{d-3} \kla\left( 1 + \frac{x^2}{3} \right).
\end{align}
Für $d<3$ existiert also ein Term $\frac{1}{x^\alpha}$ mit $\alpha \ge 1$, womit das Integral an der unteren Grenze divergiert.
\\
Für $d=3$ bekommen wir hier keine Probleme. 
Die einzige Möglichkeit für ein divergentes Integral wären dann Polstellen der Zustandsdichte.
Da aber in drei Dimensionen in der Nähe kritischer Punkte $n(x) \sim |x - x_c|^{1/2}$ gilt, treten auch hier keine Divergenzen auf.
%{\color{red}FIXME}: Konsequenzen?
Die Divergenz für $d<3$ erscheint physikalisch wenig sinnvoll, da die Auslenkungen (zumindest intuitiv) beschränkt sein sollten.
Sonst könnten sich einzelne Atome mit nicht zu vernachlässigender Wahrscheinlichkeit beliebig weit von ihrer Ruheposition entfernen.
Zumindest einmal kann das Lindemann'sche Schmelzkriterium nicht mehr sinnvoll sein, da solche Kristalle dann nie schmelzen könnten (bei Graphen scheint das mit dem ``schmelzen'' auch nicht so einfach zu sein, aber einen vergleichbaren Vorgang gibt es wohl).
\subsection*{(d) Tieftemperatur-Limes im Debye-Modell}
Im Debye-Modell ist die Dispersion streng linear, also $\omega(q) = \bar{c} q$.
Die 1. BZ wird mit einer Kugel im $\textbf{q}$-Raum ersetzt, deren Radius über $\omega_D \equiv \bar{c}q_D$ mit der Debye-Frequenz verknüpft ist.
Daraus folgt dann $n(\omega) = \frac{9}{\omega_D^3}\omega^2 \Theta(\omega_D - \omega)$ mit $\omega_D=\bar{c} \sqrt[3]{\frac{6\pi}{V_\tecxt\text{EZ}}}$.
Einsetzen in das Integral liefert
\begin{align}
\bra\langle \textbf{u}_n^2 \rangle &= \frac{\hbar}{2M} \intx\int_{0}^{\omega_D} \frac{9\omega}{\omega_D^3} \coth\kla\left( \frac{\beta\hbar\omega}{2} \right) \, \mathrm{d}\omega.
\end{align}
Für $\beta \rightarrow \infty$ vereinfacht sich das Integral zu
\begin{align}
\bra\langle \textbf{u}_n^2 \rangle &\longrightarrow \frac{9\hbar}{2M\omega_D^3} \intx\int_{0}^{\omega_D}\omega \, \mathrm{d}\omega = \frac{9\hbar}{4M\omega_D}.
\end{align}
\subsection*{(e) Hochtemperatur-Limes im Debye-Modell}
Wir substituieren $x := \beta\hbar\omega$ und erhalten
\begin{align}
\bra\langle \textbf{u}_n^2 \rangle &= \frac{9\hbar}{2M\omega_D} \intx\int_{0}^{x_D} \frac{x}{x_D^2} \coth\kla\left( \frac{x}{2} \right)  \, \mathrm{d}x
\end{align}
Für hohe Temperaturen geht $\beta$ gegen 0, wir können also den Integranden für $x\ll 1$ entwickeln
\begin{align}
x\coth\kla\left( \frac{x}{2} \right) &\sim x \frac{1 + \frac{(x/2)^2}{2}}{(x/2) + \frac{(x/2)^3}{6}} \sim 2\kla\left( 1 + \frac{(x/2)^2}{2} \right)\kla\left( 1 - \frac{(x/2)^2}{6} \right) \sim 1 + \frac{x^2}{12} .
\end{align}
Damit erhalten wir
\begin{align}
\bra\langle \textbf{u}_n^2 \rangle &\sim \frac{9k_BT}{2M\omega_D^2} \kla\left( 1 + \frac{1}{36} \kla\left( \frac{\hbar\omega_D}{k_BT} \right)^2 \right).
\end{align}
\subsection*{(f) Schmelztemperatur}
Nach dem Lindemann'schen Schmelzkriterium schmilzt ein Kristall, wenn $\bra\langle \textbf{u}_n^2 \rangle = x_Sa^2$.
Wir finden also 
\begin{align}
x_Sa^2 &= \frac{9k_BT_S}{2M\omega_D^2} \kla\left( 1 + \frac{1}{36} \kla\left( \frac{\hbar\omega_D}{k_BT_S} \right)^2 \right),\\
\Rightarrow\ \ \ 0 &= T_S^2 - \frac{2M\omega_D^2x_Sa^2}{9k_B} T_S + \frac{\hbar^2\omega_D^2}{36k_B^2}.
\end{align}
Daraus folgt dann
\begin{align}
T_S &= \frac{M\omega_D^2x_Sa^2}{9k_B} \pm \sqrt{\kla\left( \frac{M\omega_D^2x_Sa^2}{9k_B} \right)^2 - \frac{\hbar^2\omega_D^2}{36k_B^2}} = \\
&= \frac{M\omega_D^2x_Sa^2}{9k_B}\klam\left[ 1 \pm \sqrt{1 - \frac{9\hbar^2}{4\omega_D^2M^2x_S^2a^4}} \right].
\end{align}

\section{Anharmonische Korrekturen der Gitterschwingungen}
Wir betrachten einen Oszillator beschrieben durch den Hamiltonoperator $\mathcal{H}_h = \frac{p^2}{2M} + \frac{k}{2}x^2$ unter dem Einfluss einer kleinen anharmonischen Störung $V_\tecxt\text{ah} = -sx^3 - cx^4$.
\subsection*{(a) Störungsrechnung in der Quantentheorie}
Allgemein schreiben wir in der Störungstheorie die Zustände als Reihe
\begin{align}
\Ket | n \rangle &= \sum_{k\ge 0} \epsilon^k \Ket | n_k \rangle,
\end{align}
wobei üblicherweise $\epsilon\ll 1$ gelten sollte.
Der Erwartungswert eines (hermitischen) Operators $\hat{A}$ ergibt sich dann zu
\begin{align}
\bra\langle \hat{A} \rangle &= \braket\langle n | \hat{A} | n \rangle = \braket\langle n_0 | \hat{A} | n_0 \rangle + 2\epsilon \braket\langle n_0 | \hat{A} | n_1 \rangle + \mathcal{O}(\epsilon^2).
\end{align}
In unserem Fall ist
\begin{align}
\bra\langle x \rangle &= \braket\langle n_0 | x | n_0 \rangle = x_0 \braket\langle n_0 | b + b^\dagger | n_0 \rangle = 0.
\end{align}
Für die niedrigste nicht-verschwindende Ordnung ist der Term $\braket\langle n_0 | \hat{x} | n_1 \rangle$ hinreichend, da er nicht verschwindet -- das ist a priori allerdings nicht klar.
Wir benötigen also die erste Korrektur zu den Zuständen, wobei wir die Terme einzeln betrachten.
Da wir außerdem nur den Erwartungswert von $\hat{x}$ bestimmen wollen, für den $\Bra\langle n_0 |\hat{x} \propto \Bra\langle n_0+1 |\sqrt{n+1} + \Bra\langle n_0-1 |\sqrt{n}$ gilt, und für $V\propto x^4$ nur Terme der Form $\Ket | n_0\pm 2k \rangle$ für $k\in\mathbb{N}$ vorkommen, müssen wir nur den kubischen Term berücksichtigen\footnote{Da mir erst später aufgefallen ist, dass man das in dieser Ordnung gar nicht braucht, ist die Rechnung für $V\propto x^4$ weiter unten noch mit angegeben}.
Für $V\sim x^3$ finden wir 
\begin{align*}
\Ket | n_1 \rangle &= \sum_{m\neq n} \frac{\braket\langle m_0 | x^3 | n_0 \rangle}{E_{n0} - E_{m0}} \Ket | m_0 \rangle = \frac{x_0^3}{\hbar\omega} \sum_{m\neq n} \frac{\braket\langle m_0 | \kla\left( b+b^\dagger \right)^3 | n_0 \rangle}{ n - m } \Ket | m_0 \rangle = \numberthis \\
%&= \frac{x_0^3}{\hbar\omega} \sum_{m\neq n} \frac{\Ket | m_0 \rangle}{n-m} \braket\langle m_0 | b^3 + b^2b^\dagger + bb^\dagger b + b^\dagger b^2 + b^{\dagger 2}b + b^\dagger b b^\dagger + bb^{\dagger 2} + b^{\dagger 3} | n_0 \rangle = \numberthis  \\
&= \frac{x_0^3}{\hbar\omega} \sum_{m\neq n} \frac{\Ket | m_0 \rangle}{n-m} \braket\langle m_0 | b^3 + b^2b^\dagger + bb^\dagger b + b^\dagger b^2 + \tecxt\text{h.c.} | n_0 \rangle =  \numberthis \\
&= \dots \sqrt{(m+1)(m+2)(m+3)}\delta_{n,m+3} + \\
&\hspace{0.5cm} + \kla\left( \sqrt{(m+1)(m+2)(n+1)} + \sqrt{(m+1)}n + m\sqrt{n} \right) \delta_{n,m+1} + (n\leftrightarrow m) =  \numberthis \\
%&\hspace{0.5cm} + \kla\left( \sqrt{mn} + \sqrt{m(n+1)} + \sqrt{(m+1)(n+1)(n+2)} \right) \delta_{n,m-1} + \\
%&\hspace{0.5cm} + \sqrt{(n+1)(n+2)(n+3)} \delta_{n,m-3} =  \numberthis \\
&= \frac{x_0^3}{\hbar\omega} \bigg[ \sqrt{n(n-1)(n-2)} \frac{\Ket | n_0-3 \rangle}{3} - \sqrt{(n+1)(n+2)(n+3)} \frac{\Ket | n_0+3 \rangle}{3}\\
&\hspace{1.5cm} + \kla\left( \sqrt{n}(n+1) + n\sqrt{n} + (n-1)\sqrt{n} \right) \Ket | n_0-1 \rangle - \\
&\hspace{1.5cm} - \kla\left( \sqrt{n+1}(n+2) + \sqrt{n+1}(n+1) + n\sqrt{n+1} \right) \Ket | n_0+1 \rangle \bigg] =  \numberthis \\
&= \frac{x_0^3}{\hbar\omega} \bigg[ \sqrt{n(n-1)(n-2)} \frac{\Ket | n_0-3 \rangle}{3} - \sqrt{(n+1)(n+2)(n+3)} \frac{\Ket | n_0+3 \rangle}{3} + \\
&\hspace{1.5cm} + 3n\sqrt{n} \Ket | n_0 - 1 \rangle - 3(n+1)\sqrt{n+1} \Ket | n_0 + 1 \rangle \bigg]. \numberthis 
\end{align*}
Der Erwartungswert lautet dann
\begin{align}
\bra\langle \hat{x} \rangle &= 2\braket\langle n_0 | \hat{x} | -sn_1^{(x^3)} \rangle = -2sx_0\braket\langle n_0 | b + b^\dagger | n_1^{(x^3)} \rangle = \\
&= -2sx_0 \kla\left( \Bra\langle n_0+1 |\sqrt{n+1} + \Bra\langle n_0-1 |\sqrt{n} \right) \Ket | n_1^{(x^3)} \rangle = \\
&= -\frac{2sx_0^4}{\hbar\omega} \klam\left[ 3n^2 - 3(n+1)^2 \right] = \frac{12sx_0^4}{\hbar\omega} \klam\left[ n + \frac{1}{2} \right].
\end{align}
Dabei ist $x_0=\sqrt{\frac{\hbar}{2m\omega}}$, also $\bra\langle \hat{x} \rangle = \frac{3s\hbar}{m^2\omega^3} \kla\left( n+\frac{1}{2} \right)$.
\subsection*{(b) Thermodynamische Betrachtung}
Der thermodynamische Erwartungswert ergibt sich zu
\begin{align}
\bra\langle \hat{x} \rangle &= \sum_{n\ge 0} \bra\langle \hat{x} \rangle_n \e^{-\beta\hbar\omega \kla\left( n+\frac{1}{2} \right)} = \frac{3s\hbar}{m^2\omega^3} \sum_{n\ge 0} \frac{-\partial_\beta}{\hbar\omega} \e^{-\beta\hbar\omega\kla\left( n+\frac{1}{2} \right)} = \\
&= \frac{-3s\hbar}{m^2\omega^3} \frac{\partial_\beta}{\hbar\omega} \frac{1}{2\sinh\kla\left( \frac{\beta\hbar\omega}{2} \right)} = \frac{3s\hbar}{4m^2\omega^3} \frac{\cosh\kla\left( \frac{\beta\hbar\omega}{2} \right)}{\sinh\kla\left( \frac{\beta\hbar\omega}{2} \right)^2}.
\end{align}
Für $z\ll 1$ gilt
\begin{align}
\frac{\cosh(z)}{\sinh(z)^2} &\sim \frac{1 + \frac{z^2}{2}}{\kla\left( z + \frac{z^3}{6} \right)^2} \sim \frac{1}{z^2}\kla\left( 1 + \frac{z^2}{2} \right) \kla\left( 1 - \frac{z^2}{6} \right)^2 \sim\\
&\sim \frac{1}{z^2} \kla\left( 1 + \frac{z^2}{2} \right) \kla\left( 1 - \frac{z^2}{3} \right) \sim \frac{1}{z^2} \kla\left( 1 + \frac{z^2}{6} \right).
\end{align}
Im Hochtemperaturlimes $\beta\rightarrow 0$ gilt dann näherungsweise
\begin{align}
\bra\langle \hat{x} \rangle &\sim \frac{3s}{m^2\beta^2 \hbar\omega^5}  \kla\left( 1 + \frac{(\beta\hbar\omega)^2}{24} \right).
\end{align}
Um den thermodynamischen Erwartungswert zu berechnen, könnten wir zudem auch noch die Energiekorrektur in erster Ordnung berücksichtigen.
Für diese ist nur der $x^4$-Term relevant, und dort nur die Terme mit gleicher Anzahl an Erzeugern und Vernichtern.
Man findet
\begin{align}
E_{n1} &= \braket\langle n_0 | V | n_0 \rangle = -cx_0^4 \braket\langle n_0 | (b+b^\dagger)^4 | n_0 \rangle = \\
&= -cx_0^4 \braket\langle n_0 | b^2b^{\dagger 2} + bb^\dagger bb^\dagger + bb^{\dagger 2}b + b^{\dagger 2} b^2 + b^\dagger bb^\dagger b + b^\dagger b^2 b^\dagger | n_0 \rangle = \\
&= -cx_0^4 \klam\left[ (n+1)(n+2) + (n+1)^2 + (n+1)n + n(n-1) + n^2 + n(n+1) \right] = \\
&= -cx_0^4 \klam\left[ 6n^2 + 6n + 3 \right] = -6cx_0^4 \kla\left( n+\frac{1}{2} \right)^2 - \frac{3cx_0^4}{2}.
\end{align}
Damit verändert sich allerdings auch der Erwartungswert in der thermodynamischen Betrachtung, da
\begin{align}
\bra\langle \hat{x} \rangle &= \sum_{n\ge 0} \bra\langle \hat{x} \rangle_n \e^{-\beta\hbar\omega\kla\left( n+\frac{1}{2} \right) + \beta \frac{3c\hbar^2}{2m^2\omega^2} \kla\left( n+\frac{1}{2} \right)^2 + \beta \frac{3c\hbar^2}{8m^2\omega^2}}.
\end{align}
Diese Summe ist allerdings nicht mehr analytisch berechenbar.
Problematisch ist vor allem, dass die Reihe jetzt divergiert.
Das ist darauf zurückzuführen, dass die Störungstheorie nur für kleine $n$ gültig ist -- sicherlich nicht für $n\rightarrow\infty$.
Insofern bringt auch die Betrachtung des Hochtemperaturlimes $\beta\rightarrow 0$ hier keine neue Erkenntnis.
\subsection*{Anhang}
%Anwenden der Erzeuger und Vernichter nach links ergibt $\Bra\langle n_0-1 |\sqrt{n} + \Bra\langle n_0+1 |\sqrt{n+1}\dots $, weshalb nur der $n_1^{(x^3)}$-Term beiträgt.
%Wir erhalten also
%\begin{align}
%\bra\langle x \rangle &= \frac{2sx_0^4}{\hbar\omega} \klam\left[ 3n^2 - 3(n+1)^2 \right] = -\frac{24x_0^4}{\hbar\omega} \klam\left[ n + \frac{1}{2} \right].
%\end{align}
%%%%%%% BEGIN HORRIBLE CALCULATION OF THE NEXT ORDER
%In erster Ordnung Störungsrechnung ergibt sich
%\begin{align}
%\bra\langle x \rangle_1 &= \braket\langle n_0 | x | n_0 \rangle = x_0 \braket\langle n_0 | b + b^\dagger | n_0 \rangle = 0.
%\end{align}
%Wir müssen also zur zweiten Ordnung gehen.
%Dafür benötigen wir die Korrekturen zu den Zuständen, wobei wir die Terme einzeln betrachten.
%Für $V\sim x^3$ finden wir
%\begin{align*}
%\Ket | n_1 \rangle &= \sum_{m\neq n} \frac{\braket\langle m_0 | x^3 | n_0 \rangle}{E_{n0} - E_{m0}} \Ket | m_0 \rangle = \frac{x_0^3}{\hbar\omega} \sum_{m\neq n} \frac{\braket\langle m_0 | \kla\left( b+b^\dagger \right)^3 | n_0 \rangle}{ n - m } \Ket | m_0 \rangle = \numberthis \\
%%&= \frac{x_0^3}{\hbar\omega} \sum_{m\neq n} \frac{\Ket | m_0 \rangle}{n-m} \braket\langle m_0 | b^3 + b^2b^\dagger + bb^\dagger b + b^\dagger b^2 + b^{\dagger 2}b + b^\dagger b b^\dagger + bb^{\dagger 2} + b^{\dagger 3} | n_0 \rangle = \numberthis  \\
%&= \frac{x_0^3}{\hbar\omega} \sum_{m\neq n} \frac{\Ket | m_0 \rangle}{n-m} \braket\langle m_0 | b^3 + b^2b^\dagger + bb^\dagger b + b^\dagger b^2 + \tecxt\text{h.c.} | n_0 \rangle =  \numberthis \\
%&= \dots \sqrt{(m+1)(m+2)(m+3)}\delta_{n,m+3} + \\
%&\hspace{0.5cm} + \kla\left( \sqrt{(m+1)(m+2)(n+1)} + \sqrt{(m+1)}n + m\sqrt{n} \right) \delta_{n,m+1} + (n\leftrightarrow m) =  \numberthis \\
%%&\hspace{0.5cm} + \kla\left( \sqrt{mn} + \sqrt{m(n+1)} + \sqrt{(m+1)(n+1)(n+2)} \right) \delta_{n,m-1} + \\
%%&\hspace{0.5cm} + \sqrt{(n+1)(n+2)(n+3)} \delta_{n,m-3} =  \numberthis \\
%&= \frac{x_0^3}{\hbar\omega} \bigg[ \sqrt{n(n-1)(n-2)} \frac{\Ket | n_0-3 \rangle}{3} - \sqrt{(n+1)(n+2)(n+3)} \frac{\Ket | n_0+3 \rangle}{3}\\
%&\hspace{1.5cm} + \kla\left( \sqrt{n}(n+1) + n\sqrt{n} + (n-1)\sqrt{n} \right) \Ket | n_0-1 \rangle - \\
%&\hspace{1.5cm} - \kla\left( \sqrt{n+1}(n+2) + \sqrt{n+1}(n+1) + n\sqrt{n+1} \right) \Ket | n_0+1 \rangle \bigg] =  \numberthis \\
%&= \frac{x_0^3}{\hbar\omega} \bigg[ \sqrt{n(n-1)(n-2)} \frac{\Ket | n_0-3 \rangle}{3} - \sqrt{(n+1)(n+2)(n+3)} \frac{\Ket | n_0+3 \rangle}{3} + \\
%&\hspace{1.5cm} + 3n\sqrt{n} \Ket | n_0 - 1 \rangle - 3(n+1)\sqrt{n+1} \Ket | n_0 + 1 \rangle \bigg]. \numberthis 
%\end{align*}
%Für $V\sim x^4$ tragen nur die Terme mit asymmetrischer Zahl an $b$ bzw. $b^\dagger$ bei, da die Summe über $m\neq n$ läuft.
%Wir finden
%\begin{align*}
%\Ket | n_1 \rangle &= \sum_{m\neq n} \frac{\braket\langle m_0 | x^4 | n_0 \rangle}{E_{n0} - E_{m0}} \Ket | m_0 \rangle = \frac{x_0^4}{\hbar\omega} \sum_{m\neq n} \frac{\braket\langle m_0 | \kla\left( b+b^\dagger \right)^4 | n_0 \rangle}{ n - m } \Ket | m_0 \rangle = \numberthis \\
%&= \frac{x_0^4}{\hbar\omega} \sum_{m\neq n} \frac{\Ket | m_0 \rangle}{n - m} \braket\langle m_0 | b^4 + b^3b^\dagger + b^2b^\dagger b + bb^\dagger b^2 + b^\dagger b^3 + \tecxt\text{h.c.} | n_0 \rangle = \numberthis \\
%&= \dots \sqrt{(m+1)(m+2)(m+3)(m+4)} \delta_{n,m+4} + \\
%&\hspace{0.5cm} + \Big( \sqrt{(m+1)(m+2)(m+3)(n+1)} + \sqrt{(m+1)(m+2)} n + \\
%&\hspace{1.5cm} + \sqrt{(m+1)n}(n-1) + m\sqrt{n(n-1)} \Big) \delta_{n,m+2} + (n\leftrightarrow m) = \numberthis \\
%&= \frac{x_0^4}{\hbar\omega} \bigg[ \sqrt{n(n-1)(n-2)(n-3)} \frac{\Ket | n_0 - 4 \rangle}{4} - \\
%&\hspace{1.5cm} - \sqrt{(n+1)(n+2)(n+3)(n+4)} \frac{\Ket | n_0 + 4 \rangle}{4} + \\
%&\hspace{1.5cm} + \Big( \sqrt{(n-1)n}(n+1) + \sqrt{(n-1)n}n + \\
%&\hspace{2.5cm} + \sqrt{(n-1)n}(n-1) + (n-2)\sqrt{n(n-1)} \Big) \frac{\Ket | n_0-2 \rangle}{2} - \\
%&\hspace{1.5cm} - \Big( \sqrt{(n+1)(n+2)}(n+3) + \sqrt{(n+1)(n+2)}(n+2) + \\
%&\hspace{2.5cm} + \sqrt{(n+1)(n+2)}(n+1) + n\sqrt{(n+1)(n+2)} \Big) \frac{\Ket | n_0 + 2 \rangle}{2} \bigg] =  \numberthis \\
%&= \frac{x_0^4}{\hbar\omega} \bigg[ \sqrt{n(n-1)(n-2)(n-3)} \frac{\Ket | n_0 - 4 \rangle}{4} - \\
%&\hspace{1.5cm} - \sqrt{(n+1)(n+2)(n+3)(n+4)} \frac{\Ket | n_0 + 4 \rangle}{4} + \\
%&\hspace{1.5cm} + \sqrt{n(n-1)}(2n-1)\Ket | n_0 - 2 \rangle - \sqrt{(n+1)(n+2)} (2n+3)\Ket | n_0 + 2 \rangle \bigg]. \numberthis
%\end{align*}
%Insgesamt erhalten wir dann die Korrektur
%\begin{align}
%\bra\langle x \rangle_2 &= \braket\langle n_0 + n_1 | x | n_0 + n_1 \rangle = 2\braket\langle n_0 | x | n_1 \rangle + \braket\langle n_1 | x | n_1 \rangle = \\
%&= 2x_0 \braket\langle n_0 | b + b^\dagger | sn_1^{(x^3)} + cn_1^{(x^4)} \rangle + x_0\braket\langle sn_1^{(x^3)} + cn_1^{(x^4)} | b + b^\dagger | sn_1^{(x^3)} + cn_1^{(x^4)} \rangle.
%\end{align}
%Im ersten Teil haben wir $\Bra\langle n_0-1 |\sqrt{n} + \Bra\langle n_0+1 |\sqrt{n+1}\dots $, weshalb nur der $n_1^{(x^3)}$-Teil beiträgt.
%Hier bekommen wir also
%\begin{align}
%\bra\langle x \rangle_2 &= \frac{2x_0^4}{\hbar\omega} \klam\left[ 3n^2 - 3(n+1)^2 \right] = -\frac{6x_0^4}{\hbar\omega} \klam\left[ 2n + 1 \right].
%\end{align}
%Auch hier tragen nur die gemischten Terme bei, also
%\begin{align*}
%\bra\langle x \rangle_2 &= 2x_0sc \braket\langle n_1^{(x^3)} | b + b^\dagger | n_1^{(x^4)} \rangle = \numberthis \\
%&= \frac{2x_0^8sc}{\hbar^2\omega^2} \bigg[ \sqrt{n(n-1)(n-2)} \frac{\Bra\langle n_0 - 3 |}{3} - \sqrt{(n+1)(n+2)(n+3)} \frac{\Bra\langle n_0 + 3 |}{3} + \\
%&\hspace{1.5cm} + 3n\sqrt{n} \Bra\langle n_0 - 1 | - 3(n+1)\sqrt{n+1} \Bra\langle n_0 + 1 | \bigg] (b+b^\dagger) \\
%&\hspace{1.5cm}\bigg[ \sqrt{n(n-1)(n-2)(n-3)} \frac{\Ket | n_0 - 4 \rangle}{4} - \\
%&\hspace{1.5cm} - \sqrt{(n+1)(n+2)(n+3)(n+4)} \frac{\Ket | n_0 + 4 \rangle}{4} + \\
%&\hspace{1.5cm} + \sqrt{n(n-1)}(2n-1)\Ket | n_0 - 2 \rangle - \sqrt{(n+1)(n+2)} (2n+3)\Ket | n_0 + 2 \rangle \bigg] =  \numberthis \\
%&= \frac{2x_0^8sc}{\hbar^2\omega^2} \bigg[  \frac{n(n-1)(n-2)}{3} (2n-1) +  \frac{n(n-1)(n-2)(n-3)}{12} + \\
%&\hspace{1.5cm} + \frac{(n+1)(n+2)(n+3)(n+4)}{12} +  \frac{(n+1)(n+2)(n+3)}{3} (2n+3) + \\
%&\hspace{1.5cm} + 3n^2(n-1) (2n-1) + 3(n+1)(n+1)(n+2) (2n+3) \bigg] =  \numberthis \\
%&= \frac{2x_0^8sc}{\hbar^2\omega^2} \bigg[ n^4\kla\left( \frac{4}{3} + \frac{2}{12} + 12 \right) + n^3\kla\left( -\frac{7}{3} - \frac{6}{12} + \frac{10}{12} + \frac{15}{3} - 9 + 33 \right) + \\
%&\hspace{1.5cm} + n^2 \kla\left( \frac{7}{3} + \frac{11}{12} + \frac{35}{12} + \frac{40}{3} + 3 + 22 \right) + n\kla\left( \frac{-2}{3} - \frac{6}{12} + \frac{50}{12} + \frac{39}{3} + 57 \right) + \\
%&\hspace{1.5cm} + \frac{24}{12} + \frac{18}{3} + 18 \bigg] = \numberthis \\
%&= \frac{2x_0^8sc}{\hbar^2\omega^2} \bigg[ \frac{27n^4}{2} + 21n^3 + \frac{58n^2}{3} + 73n + 26 \bigg]  \numberthis 
%\end{align*}
%Insgesamt erhalten wir also
%\begin{align}
%\bra\langle x \rangle_2 &= -\frac{6x_0^4}{\hbar\omega} \klam\left[ 2n + 1 \right] + \frac{x_0^8sc}{\hbar^2\omega^2} \bigg[ 27n^4 + 42n^3 + \frac{116n^2}{3} + 146n + 52 \bigg].
%\end{align}
%%%%%%% END HORRIBLE CALCULATION OF THE NEXT ORDER
Hier noch die Rechnung für $V\sim x^4$.
Es tragen nur die Terme mit asymmetrischer Zahl an $b$ bzw. $b^\dagger$ bei, da die Summe über $m\neq n$ läuft.
Wir finden
\begin{align*}
\Ket | n_1 \rangle &= \sum_{m\neq n} \frac{\braket\langle m_0 | x^4 | n_0 \rangle}{E_{n0} - E_{m0}} \Ket | m_0 \rangle = \frac{x_0^4}{\hbar\omega} \sum_{m\neq n} \frac{\braket\langle m_0 | \kla\left( b+b^\dagger \right)^4 | n_0 \rangle}{ n - m } \Ket | m_0 \rangle = \numberthis \\
&= \frac{x_0^4}{\hbar\omega} \sum_{m\neq n} \frac{\Ket | m_0 \rangle}{n - m} \braket\langle m_0 | b^4 + b^3b^\dagger + b^2b^\dagger b + bb^\dagger b^2 + b^\dagger b^3 + \tecxt\text{h.c.} | n_0 \rangle = \numberthis \\
&= \dots \sqrt{(m+1)(m+2)(m+3)(m+4)} \delta_{n,m+4} + \\
&\hspace{0.5cm} + \Big( \sqrt{(m+1)(m+2)(m+3)(n+1)} + \sqrt{(m+1)(m+2)} n + \\
&\hspace{1.5cm} + \sqrt{(m+1)n}(n-1) + m\sqrt{n(n-1)} \Big) \delta_{n,m+2} + (n\leftrightarrow m) = \numberthis \\
&= \frac{x_0^4}{\hbar\omega} \bigg[ \sqrt{n(n-1)(n-2)(n-3)} \frac{\Ket | n_0 - 4 \rangle}{4} - \\
&\hspace{1.5cm} - \sqrt{(n+1)(n+2)(n+3)(n+4)} \frac{\Ket | n_0 + 4 \rangle}{4} + \\
&\hspace{1.5cm} + \Big( \sqrt{(n-1)n}(n+1) + \sqrt{(n-1)n}n + \\
&\hspace{2.5cm} + \sqrt{(n-1)n}(n-1) + (n-2)\sqrt{n(n-1)} \Big) \frac{\Ket | n_0-2 \rangle}{2} - \\
&\hspace{1.5cm} - \Big( \sqrt{(n+1)(n+2)}(n+3) + \sqrt{(n+1)(n+2)}(n+2) + \\
&\hspace{2.5cm} + \sqrt{(n+1)(n+2)}(n+1) + n\sqrt{(n+1)(n+2)} \Big) \frac{\Ket | n_0 + 2 \rangle}{2} \bigg] =  \numberthis \\
&= \frac{x_0^4}{\hbar\omega} \bigg[ \sqrt{n(n-1)(n-2)(n-3)} \frac{\Ket | n_0 - 4 \rangle}{4} - \\
&\hspace{1.5cm} - \sqrt{(n+1)(n+2)(n+3)(n+4)} \frac{\Ket | n_0 + 4 \rangle}{4} + \\
&\hspace{1.5cm} + \sqrt{n(n-1)}(2n-1)\Ket | n_0 - 2 \rangle - \sqrt{(n+1)(n+2)} (2n+3)\Ket | n_0 + 2 \rangle \bigg]. \numberthis
\end{align*}

\end{document}